\section{Introduction}
\label{sec:intro}

The anomalous magnetic moment of the muon, \(a_\mu = (g-2)/2\), is a cornerstone quantity for testing the Standard Model (SM) and probing new physics. In 2025, the Muon \(g-2\) Experiment at Fermilab released its final measurement~\cite{Aguillard2025,Driutti2025} based on the complete Run 1--6 dataset, yielding a world average experimental value combining with the BNL E821 result~\cite{Bennett2006} of
\[
a_\mu^{\mathrm{exp}} = 116\,592\,072(1.5) \times 10^{-11} \quad (124\ \mathrm{ppb}).
\]

Within the Standard Model, \(a_\mu^{\mathrm{SM}}\) is decomposed into several contributions:
\[
a_\mu^{\mathrm{SM}} = a_\mu^{\mathrm{QED}} + a_\mu^{\mathrm{EW}} + a_\mu^{\mathrm{HVP}} + a_\mu^{\mathrm{HLbL}},
\]
where \(a_\mu^{\mathrm{QED}}\) denotes the quantum electrodynamics contribution (computed to five loops), \(a_\mu^{\mathrm{EW}}\) the electroweak contribution, \(a_\mu^{\mathrm{HVP}}\) the hadronic vacuum polarization, and \(a_\mu^{\mathrm{HLbL}}\) the hadronic light-by-light scattering. Among these, the hadronic contributions --- particularly \(a_\mu^{\mathrm{HVP}}\) --- dominate the theoretical uncertainty.

The largest uncertainty in the SM prediction comes from the leading-order hadronic vacuum polarization (LO-HVP) contribution, \(a_\mu^{\mathrm{HVP}}\), which cannot be computed perturbatively. In the traditional data-driven dispersive approach, \(a_\mu^{\mathrm{HVP}}\) is obtained via the dispersion integral~\cite{Brodsky1968,Lautrup1968}
\[
a_\mu^{\mathrm{HVP}} = \frac{\alpha^2}{3\pi^2} \int_{4m_\pi^2}^{\infty} \frac{ds}{s} K(s) \, R(s),
\]
where \(R(s) = \sigma(e^+e^-\to \mathrm{hadrons}) / \sigma(e^+e^-\to \mu^+\mu^-)\) is the hadronic \(R\)-ratio, and \(K(s)\) is a known QED kernel function that peaks at low energies. Consequently, high-precision measurements of exclusive hadronic cross sections at \(e^+e^-\) colliders are indispensable for a reliable SM prediction.

The WP2020~\cite{WP2020} evaluation by the Muon \(g-2\) Theory Initiative relied on this data-driven method. It yielded a SM prediction of
\[
a_\mu^{\mathrm{SM}}(\mathrm{WP2020}) = 116\,591\,810(43) \times 10^{-11},
\]
which, when compared with the then-available experimental value, resulted in a \(>5\sigma\) deviation, fueling intense interest in possible new physics.

However, the theoretical landscape has undergone a fundamental shift with the publication of the 2025 White Paper (WP2025)~\cite{WP2025}. Persistent tensions among different \(e^+e^-\to\pi^+\pi^-\) measurements --- in particular the higher \(\pi\pi\) cross section reported by the CMD-3 experiment~\cite{CMD3_Ignatov2024} compared to other earlier experiments --- have made it impossible to produce a meaningful combined data-driven average. In contrast, lattice QCD~\cite{Borsanyi2021} calculations of the LO-HVP have reached sub-percent precision and show good internal consistency among different collaborations.

Consequently, WP2025 adopts a consolidated lattice-QCD average as its central value for the LO-HVP contribution, rather than the data-driven result used in WP2020. This change shifts the total SM prediction upward to
\[
a_\mu^{\mathrm{SM}}(\mathrm{WP2025}) = 116\,592\,033(62) \times 10^{-11}.
\]

The comparison between theory and experiment now yields:
\[
a_\mu^{\mathrm{exp}} - a_\mu^{\mathrm{SM}}(\mathrm{WP2025}) = 38(63) \times 10^{-11},
\]
corresponding to a negligible deviation of approximately \(0.6\sigma\). The long-standing \(>5\sigma\) discrepancy has effectively been resolved with the updated SM prediction. Nevertheless, this resolution does not imply that the problem is fully understood. The disagreement between the data-driven and lattice-QCD evaluations of the LO-HVP term remains unresolved, as does the tension between CMD-3 and other \(e^+e^-\to\pi^+\pi^-\) experiments. Resolving these discrepancies requires new, independent, high-precision measurements of the key hadronic cross sections, particularly those with the largest contributions to \(a_\mu^{\mathrm{HVP}}\).

% Among all exclusive channels, the \(2\pi\) final state dominates, contributing \(\sim 70\%\) of the total HVP. The second most important channel is \(e^+e^-\to\pi^+\pi^-\pi^0\) (\(3\pi\)), which contributes \(\sim 10\%\) but is a major source of systematic uncertainty due to its complex resonance structure (dominated by \(\omega\) and \(\phi\)). 

For the \(3\pi\) channel, the Belle II experiment reported a new measurement in 2024~\cite{BelleII_3pi_2024}. However, concerns have been raised regarding the reliability of its uncertainty estimation, leading some recent SM theory compilations to explicitly exclude Belle II data from their averages\cite{Hoferichter2025}. This situation highlights the urgent need for independent cross section measurements from other facilities.

The BESIII experiment, with its large data samples, is ideally suited to provide a reliable and independent measurement of the \(e^+e^- \to \pi^+\pi^-\pi^0\) cross section in the low-energy region critical for \(a_\mu^{\mathrm{HVP}}\). In this memo, we use an \(e^+e^-\) collision sample corresponding to an integrated luminosity of 20~fb\(^{-1}\) collected at a center-of-mass energy of 3.773~GeV. We present a measurement based on the initial-state radiation (ISR) technique using the process \(e^+e^- \to \pi^+\pi^-\pi^0 \gamma_{\mathrm{ISR}}\), which allows access to the \(3\pi\) hadronic cross section from continuum production without energy scanning. This work aims to provide a valid and systematically well-controlled measurement, thereby adding reliable experimental input for the \(3\pi\) channel and contributing to the ongoing effort to reconcile the data-driven and lattice-QCD approaches.